\appendixpage
\section{Instance parameters and resource limits}
\label{app:tiers}
We use three parameter ranges: A1 contains small familiar instances, while A2 and A3 increase size. Family selection for A3 screened the first sampled setting for a ratio of at least 1.4 between the implemented baseline $z$ and a bound or target $b$: $b/z$ for maximisation and $z/b$ for minimisation. This screening identifies families with room for improvement; the remaining instances vary the parameters within those families. Answers must also fit the representation limits and be verifiable within 60 seconds. Sorting networks and Kakeya are evaluated separately as controls. For example, the Kakeya setting $(q,n)=(13,3)$ has a proven lower bound of 595~\citep{bukh2021kakeya} and a construction of size 703, leaving a ratio of only 1.18.

The main evaluation contains 69 instances: 61 in the twelve non-code families, four Shannon instances and four trifference instances. Each configuration contributes one response per instance under the protocol in Appendix~\ref{app:runner}. The released manifest specifies the selected instances. Table~\ref{tab:tiers} lists the instances, and Table~\ref{tab:triangle_domains} specifies the triangular domains. Kakeya provides a separate control for the family-selection comparison. Smaller-instance parameters appear in Table~\ref{tab:aux_parameters}.
\small\setlength{\tabcolsep}{4pt}
\begin{longtable}{@{}P{.30\linewidth}P{.67\linewidth}@{}}
\caption{The 69 instances in the main evaluation. Each configuration contributes one response per instance.}\label{tab:tiers}\\
\toprule
Family & Parameters \\
\midrule\endfirsthead
\toprule
Family & Parameters \\
\midrule\endhead
\bottomrule\endfoot
Cap sets & $d=10$ \\
Cap sets & $d=11$ \\
Cap sets & $d=9$ \\
\addlinespace[4pt]
Corners & $n=138$ \\
Corners & $n=183$ \\
Corners & $n=150$ \\
Corners & $n=162$ \\
Corners & $n=185$ \\
\addlinespace[4pt]
Schur & $k=7$ \\
Schur & $k=6$ \\
Schur & $k=8$ \\
\addlinespace[4pt]
$\mathbb F_q^n$ AP-free & $n=6,\ q=5$ \\
$\mathbb F_q^n$ AP-free & $n=5,\ q=5$ \\
$\mathbb F_q^n$ AP-free & $n=4,\ q=7$ \\
\addlinespace[4pt]
Spherical ($60^\circ$) & $d=14$ \\
Spherical ($60^\circ$) & $d=15$ \\
Spherical (variable) & $c=24/100,\ d=12$ \\
Spherical (variable) & $c=39/100,\ d=12$ \\
Spherical (variable) & $c=36/100,\ d=14$ \\
Spherical (variable) & $c=39/100,\ d=14$ \\
Spherical (variable) & $c=36/100,\ d=12$ \\
Spherical ($60^\circ$) & $d=13$ \\
\addlinespace[4pt]
LABS & $N=280$ \\
LABS & $N=375$ \\
LABS & $N=357$ \\
LABS & $N=281$ \\
LABS & $N=240$ \\
\addlinespace[4pt]
Heilbronn (square) & $n=46$ \\
Heilbronn (square) & $n=26$ \\
Heilbronn (square) & $n=40$ \\
Heilbronn (square) & $n=50$ \\
Heilbronn (square) & $n=48$ \\
Heilbronn (triangle) & $n=32$, $T_{1}$ \\
Heilbronn (triangle) & $n=34$, $T_{2}$ \\
Heilbronn (triangle) & $n=31$, $T_{3}$ \\
Heilbronn (triangle) & $n=40$, $T_{4}$ \\
Heilbronn (triangle) & $n=26$, $T_{5}$ \\
\addlinespace[4pt]
Degree--diameter & $d=4,\ k=5$ \\
Degree--diameter & $d=7,\ k=3$ \\
Degree--diameter & $d=4,\ k=4$ \\
Degree--diameter & $d=3,\ k=7$ \\
Degree--diameter & $d=6,\ k=3$ \\
Degree--diameter & $d=5,\ k=4$ \\
Degree--diameter & $d=3,\ k=6$ \\
\addlinespace[4pt]
Linear-equation-free & $a=3,\ b=4,\ n=19044$ \\
Linear-equation-free & $a=2,\ b=5,\ n=18883$ \\
Linear-equation-free & $a=2,\ b=3,\ n=17965$ \\
Linear-equation-free & $a=3,\ b=4,\ n=13312$ \\
Linear-equation-free & $a=3,\ b=4,\ n=6248$ \\
\addlinespace[4pt]
Covering designs & $k=7,\ t=4,\ v=18$ \\
Covering designs & $k=6,\ t=3,\ v=30$ \\
Covering designs & $k=6,\ t=3,\ v=24$ \\
\addlinespace[4pt]
Matrix multiplication & $m=4,\ n=4,\ p=5$ \\
Matrix multiplication & $m=5,\ n=5,\ p=5$ \\
Matrix multiplication & $m=5,\ n=4,\ p=5$ \\
\addlinespace[4pt]
MOLS & $n=14$ \\
MOLS & $n=10$ \\
MOLS & $n=18$ \\
MOLS & $n=15$ \\
MOLS & $n=12$ \\
MOLS & $n=20$ \\
\addlinespace[4pt]
Shannon & $d=24,\ q=7$ \\
Shannon & $d=40,\ q=7$ \\
Shannon & $d=48,\ q=9$ \\
Shannon & $d=64,\ q=9$ \\
\addlinespace[4pt]
Trifference & $m=1,\ n=64$ \\
Trifference & $m=1,\ n=96$ \\
Trifference & $m=1,\ n=144$ \\
Trifference & $m=1,\ n=192$ \\
\addlinespace[4pt]
\end{longtable}

\normalsize
\subsection{Triangular domains}
\small\begin{longtable}{@{}lP{.27\linewidth}P{.27\linewidth}P{.27\linewidth}@{}}
\caption{Triangular domains used in the main evaluation. Coordinates are integers.}\label{tab:triangle_domains}\\
\toprule
Domain & Vertex 1 & Vertex 2 & Vertex 3 \\
\midrule\endfirsthead
\toprule
Domain & Vertex 1 & Vertex 2 & Vertex 3 \\
\midrule\endhead
\bottomrule\endfoot
$T_{1}$ & $(9026,8197)$ & $(1769,9178)$ & $(7298,1276)$ \\
$T_{2}$ & $(578,9933)$ & $(8598,7337)$ & $(2075,588)$ \\
$T_{3}$ & $(411,7999)$ & $(9004,9890)$ & $(3031,713)$ \\
$T_{4}$ & $(7922,8793)$ & $(9917,1548)$ & $(1510,3942)$ \\
$T_{5}$ & $(526,8526)$ & $(8696,8279)$ & $(8084,1688)$ \\
\end{longtable}
\normalsize
Each triangle is specified by its three integer-coordinate vertices. The objective is the minimum triangle area divided by the area of the containing domain.
\appendixpage
\subsection{Smaller-instance parameter draws}
\small\begin{longtable}{@{}P{.30\linewidth}lP{.56\linewidth}@{}}
\caption{Three representative smaller-instance draws per tier, in seed order 0, 1, 2. Trifference uses the lengths listed in the main evaluation table.}\label{tab:aux_parameters}\\
\toprule
Family & Tier & Representative parameters \\
\midrule\endfirsthead
\toprule
Family & Tier & Representative parameters \\
\midrule\endhead
\bottomrule\endfoot
Cap sets & A1 & $d=3$; $d=3$; $d=4$ \\
 & A2 & $d=8$; $d=7$; $d=7$ \\
\addlinespace[4pt]
Corner-free sets & A1 & $n=23$; $n=15$; $n=22$ \\
 & A2 & $n=66$; $n=50$; $n=79$ \\
\addlinespace[4pt]
Schur colourings & A1 & $k=4$; $k=4$; $k=3$ \\
 & A2 & $k=5$; $k=5$; $k=6$ \\
\addlinespace[4pt]
Progression-free sets over finite fields & A1 & $q=5,\ n=4$; $q=5,\ n=4$; $q=5,\ n=4$ \\
 & A2 & $q=7,\ n=3$; $q=7,\ n=4$; $q=7,\ n=4$ \\
\addlinespace[4pt]
Spherical codes (60$^\circ$) & A1 & $d=7$; $d=7$; $d=7$ \\
 & A2 & $d=11$; $d=12$; $d=9$ \\
\addlinespace[4pt]
Spherical codes (variable angle) & A1 & $d=6,\ c=37/100$; $d=6,\ c=35/100$; $d=6,\ c=48/100$ \\
 & A2 & $d=10,\ c=47/100$; $d=10,\ c=45/100$; $d=9,\ c=20/100$ \\
\addlinespace[4pt]
Low-autocorrelation binary sequences & A1 & $N=19$; $N=25$; $N=29$ \\
 & A2 & $N=68$; $N=74$; $N=91$ \\
\addlinespace[4pt]
Heilbronn triangles (square) & A1 & $n=7$; $n=6$; $n=5$ \\
 & A2 & $n=18$; $n=17$; $n=18$ \\
\addlinespace[4pt]
Heilbronn triangles (triangle) & A1 & $n=9$; $n=12$; $n=11$ \\
 & A2 & $n=16$; $n=14$; $n=12$ \\
\addlinespace[4pt]
Degree--diameter graphs & A1 & $d=4,\ k=3$; $d=3,\ k=3$; $d=3,\ k=3$ \\
 & A2 & $d=10,\ k=2$; $d=5,\ k=3$; $d=3,\ k=5$ \\
\addlinespace[4pt]
Linear-equation-free sets & A1 & $a=3,\ b=4,\ n=129$; $a=1,\ b=2,\ n=101$; $a=2,\ b=3,\ n=155$ \\
 & A2 & $a=3,\ b=5,\ n=1394$; $a=1,\ b=3,\ n=1445$; $a=3,\ b=4,\ n=1439$ \\
\addlinespace[4pt]
Covering designs & A1 & $v=12,\ k=6,\ t=3$; $v=13,\ k=4,\ t=2$; $v=10,\ k=5,\ t=3$ \\
 & A2 & $v=16,\ k=6,\ t=3$; $v=15,\ k=7,\ t=4$; $v=16,\ k=6,\ t=3$ \\
\addlinespace[4pt]
Matrix multiplication & A1 & $n=2,\ m=2,\ p=4$; $n=3,\ m=3,\ p=3$; $n=2,\ m=2,\ p=4$ \\
 & A2 & $n=3,\ m=3,\ p=5$; $n=3,\ m=3,\ p=4$; $n=4,\ m=4,\ p=5$ \\
\addlinespace[4pt]
Mutually orthogonal Latin squares & A1 & $n=12$; $n=10$; $n=10$ \\
 & A2 & $n=15$; $n=14$; $n=14$ \\
\addlinespace[4pt]
Shannon codes & A1 & $q=7,\ d=3$; $q=9,\ d=3$; $q=7,\ d=4$ \\
 & A2 & $q=9,\ d=12$; $q=7,\ d=8$; $q=7,\ d=12$ \\
\addlinespace[4pt]
\end{longtable}
\normalsize
